\chapter{Introducción}

En fabricación aeronáutica, todo movimiento de material exige registrar kit, estación, hora, operador y, en su caso, incidencias FOD. Un AMR industrial puede asumir una parte de ese trabajo: transporte repetitivo, confirmación de entregas y rutas flexibles sin carriles físicos.

Lunabotics prepara una oferta para Alestis Puerto Real. El cliente trabaja con componentes estructurales del A320, entre ellos el HTP, elemento aerodinámico horizontal situado en la cola que contribuye a la estabilidad longitudinal de la aeronave. En torno a la sección 19.1 del avión convergen operaciones de estructura posterior, interfaces con elementos de cola, útiles de montaje, inspecciones y flujos de materiales que requieren coordinación fina entre estaciones.

\section{Perfil de Lunabotics}

Lunabotics se plantea como una integradora de 24 personas: cinco roboticistas con experiencia en ROS 2/Nav2, tres ingenieros de automatización, un responsable de seguridad de máquinas, tres desarrolladores de HMI/datos, cuatro técnicos de puesta en marcha, dos ingenieros de soporte, tres perfiles de proyecto y documentación, dos perfiles comerciales y un responsable de calidad. La empresa ha desplegado AMR en logística interna de componentes aeronáuticos. Esta oferta se limita a logística auxiliar y despliegue por fases, fuera de procesos certificados de montaje estructural.

\begin{table}[H]
\centering
\begin{tabularx}{\textwidth}{P{0.23\textwidth}P{0.35\textwidth}X}
\toprule
Referencia ficticia & Trabajo realizado & Utilidad para Alestis\\
\midrule
Planta aeronáutica auxiliar & Flota de 3 AMR para suministro de consumibles y kits en zona de pretratamiento de composites. & Experiencia directa en logística de planta aeronáutica con disciplina FOD.\\
Montaje automotriz & Integración de 4 AMR tipo MiR para suministro de subconjuntos. & Gestión de misiones, seguridad en pasillos compartidos y aceptación de operarios.\\
Celda universitaria mecanum & Validación de cinemática holonómica con KUKA YouBot en CoppeliaSim. & Base técnica para maniobra lateral, control de ruedas mecanum y pruebas previas al piloto.\\
Almacén electrónico & Registro QR/RFID de entregas y dashboard de incidencias. & Trazabilidad de kits, herramientas y consumibles sin depender de memoria verbal de turno.\\
\bottomrule
\end{tabularx}
\caption{Perfil operativo de Lunabotics para la oferta.}
\label{tab:perfil-lunabotics}
\end{table}

Cuando un operario debe abandonar su estación para buscar herramientas, consumibles o documentación, el tiempo perdido no siempre aparece como parada formal, pero reduce capacidad real. Algo similar ocurre con kits incompletos, útiles no localizados, entregas sin registrar o inspecciones FOD tardías.


\section{Alcance de la propuesta}

El alcance cubre transporte interno de kits, herramientas, consumibles y útiles ligeros o moderados. También incluye inspección preventiva FOD con cámara, trazabilidad RFID/QR, retorno de útiles usados, supervisión de misiones y registro de eventos. Se excluyen la fabricación directa de piezas aeronáuticas, la manipulación estructural pesada del HTP y la sustitución de procesos certificados de calidad. La solución queda dentro de una aplicación propia de Sistemas Robóticos Móviles \cite{siegwart2011,tzafestas2014,ollero2001}.

\section{Supuestos y exclusiones}

El caso asume una fábrica interior con pavimento regular, cobertura WiFi industrial ampliable, estaciones con puntos de entrega identificables y personal disponible para carga/descarga del AMR. Los kits y herramientas se etiquetarían con RFID, QR o una combinación de ambos.

Del presupuesto quedan fuera la obra civil, modificaciones estructurales de estaciones e integraciones ERP/MES altamente personalizadas no descritas en el alcance. Con esas exclusiones, el coste se concentra en robótica móvil, kitting, sensores, registro operativo, seguridad y soporte.
